\begin{frame}{Decisiones empresariales}
\begin{itemize}
  \item Fortalecer cobranza mediante seguimiento por antigüedad de saldos.
  \item Controlar inventarios para reducir capital inmovilizado.
  \item Revisar costo financiero y calendario de deuda.
  \item Evaluar margen por producto o línea de negocio.
  \item Proyectar flujo de efectivo mensual.
  \item Monitorear indicadores con cierre contable confiable.
\end{itemize}
\end{frame}

\begin{frame}{Aplicación en Ingeniería de Sistemas}
\begin{center}
\begin{tikzpicture}[node distance=0.65cm, every node/.style={draw, rounded corners, align=center, minimum width=1.7cm, minimum height=0.7cm, font=\scriptsize}]
\node (erp) {ERP};
\node[right=of erp] (bd) {Base de\\datos};
\node[right=of bd] (calc) {Cálculo\\ratios};
\node[right=of calc] (dash) {Dashboard};
\node[right=of dash] (alert) {Alertas};
\node[right=of alert] (dec) {Decisión};
\foreach \a/\b in {erp/bd,bd/calc,calc/dash,dash/alert,alert/dec}{\draw[-{Stealth}] (\a) -- (\b);}
\end{tikzpicture}
\end{center}
\begin{itemize}
  \item SQL obtiene saldos contables.
  \item Python o Excel calcula ratios.
  \item Power BI presenta indicadores.
  \item La trazabilidad contable debe mantenerse.
\end{itemize}
\end{frame}
